\section{Incertidumbre y predicción selectiva}
\label{sec:marco_incertidumbre}

Un clasificador devuelve siempre una categoría y una probabilidad por clase, aun ante una imagen que no sabe
distinguir. En medicina una respuesta equivocada presentada con la misma seguridad que una correcta puede
llevar a una conducta errónea, y reconocer cuándo no conviene responder es un problema con teoría propia.

\subsection{Incertidumbre en la clasificación diagnóstica}
\label{sec:marco_incert_fuentes}

La incertidumbre de una predicción tiene dos fuentes. La incertidumbre aleatoria proviene de los datos: una
imagen desenfocada, o una lesión cuyo aspecto es compatible con varias categorías, deja la respuesta
indeterminada aunque el modelo sea bueno, y reunir más datos no la reduce. La incertidumbre epistémica
proviene del modelo, que no vio suficientes ejemplos de cierto tipo de lesión, y disminuye con más
datos~\cite{kendall2017uncertainties}.

Una tercera fuente es el cambio de distribución (Sección~\ref{sec:marco_generalizacion}): cuando la imagen
proviene de un equipo, centro o población que el entrenamiento no cubrió, lo aprendido puede no aplicarse.
Como la softmax de la Ecuación~\ref{eq:marco_softmax} hace que las probabilidades sumen uno, no existe una
salida que exprese «este caso no lo conozco», y la probabilidad no refleja necesariamente la incertidumbre.

En dermatología los errores no cuestan lo mismo: no detectar una lesión maligna tiene consecuencias
distintas de derivar una benigna. Además, la decisión final la toma una persona. En un experimento con
dermatólogos y médicos de atención primaria, el apoyo de un sistema de aprendizaje profundo mejoró la
exactitud de los dos grupos~\cite{groh2024learningaided}, de modo que su salida es un insumo de la decisión
clínica y no la decisión misma. Un sistema que reconoce su límite y deriva puede ser preferible a uno que
siempre responde, y formalizar la abstención permite medir cuánto se gana.

\subsection{Predicción selectiva y abstención}
\label{sec:marco_pred_selectiva}

Un clasificador con opción de rechazo se describe con un par $(f, g)$. La función $f$ asigna una categoría a
cada imagen $x$, y la función $g$ decide si la predicción se acepta ($g(x) = 1$) o si el modelo se abstiene
($g(x) = 0$). La cobertura es la fracción de imágenes que el modelo acepta, y el riesgo selectivo es la tasa
de error medida solo sobre las aceptadas,
\begin{equation}
    \mathrm{cobertura}(g) = \mathbb{E}\bigl[g(x)\bigr], \qquad
    \mathrm{riesgo}(f, g) = \frac{\mathbb{E}\bigl[\mathbf{1}[f(x) \neq y]\; g(x)\bigr]}{\mathbb{E}\bigl[g(x)\bigr]},
    \label{eq:marco_riesgo_selectivo}
\end{equation}
donde $y$ es la categoría real. Sobre $n$ imágenes, la cobertura es la proporción de imágenes con
$g(x_i) = 1$ y el riesgo, la proporción de errores entre ellas~\cite{elyaniv2010foundations,geifman2017selective}.

Una forma habitual de construir $g$ es partir de una señal de confianza $\kappa(x)$, un número que crece con la
seguridad del modelo, y de un umbral $\tau$,
\begin{equation}
    g(x) = \mathbf{1}\bigl[\kappa(x) \ge \tau\bigr].
    \label{eq:marco_umbral_confianza}
\end{equation}
Subir $\tau$ reduce la cobertura, porque se aceptan menos imágenes, y reduce también el riesgo si la señal es
informativa, porque las imágenes descartadas son las más propensas al error. El compromiso se resume en la
curva riesgo--cobertura, que dibuja el riesgo selectivo contra la cobertura al recorrer todos los umbrales.
Una señal es mejor cuanto más rápido baja el riesgo al reducir la cobertura, y una que no distingue los
aciertos de los errores deja el riesgo en el nivel del error global. Todo depende entonces de qué señal se use.

\subsection{Señales de confianza y calibración}
\label{sec:marco_senales_confianza}

La señal más simple es la máxima probabilidad de la softmax, $\kappa(x) = \max_c p_c$, que sirve de línea
base para detectar predicciones erróneas~\cite{hendrycks2017baseline}. Para que un umbral sobre ella tenga una
lectura probabilística, el modelo tiene que estar calibrado, es decir, sus probabilidades tienen que coincidir con sus frecuencias de acierto: entre las
predicciones a las que asigna confianza $p$, acierta una fracción $p$. En la notación de
\textcite{guo2017calibration}, con $\hat{Y}$ la clase predicha y $\hat{P}$ su probabilidad,
$P(\hat{Y} = Y \mid \hat{P} = p) = p$. Las redes profundas actuales tienden a ser sobreconfiadas, y las más
profundas y anchas suelen estar peor calibradas que las pequeñas aunque clasifiquen
mejor~\cite{guo2017calibration}. Con una confianza calibrada, un umbral de 0.9 acepta predicciones con al
menos 0.9 de probabilidad de ser correctas; sin calibración el umbral no tiene esa lectura. Las métricas de
calibración y sus correcciones se presentan en la Sección~\ref{sec:marco_eval_calib}.

Hay señales que no dependen de la probabilidad de un único modelo. El desacuerdo entre modelos aproxima la
incertidumbre del sistema, un principio que explotan los ensambles. Un ensamble promedia las probabilidades de
varios modelos entrenados por separado, lo que se llama voto suave, y suele ser más exacto y estar mejor
calibrado que cada uno de ellos~\cite{lakshminarayanan2017simple}.

\phantomsection\label{sec:marco_concordancia}%
Dos predictores con sesgos distintos, por ejemplo un clasificador supervisado y uno basado en vecinos
(Sección~\ref{sec:marco_vlm_clip}), pueden coincidir o no sobre la misma imagen. La \textbf{concordancia} entre
sus predicciones $\hat{y}_1(x)$ y $\hat{y}_2(x)$ es
\begin{equation}
    a(x) = \mathbf{1}\bigl[\hat{y}_1(x) = \hat{y}_2(x)\bigr],
    \label{eq:marco_concordancia}
\end{equation}
y sirve como señal solo si los casos en que los predictores coinciden se equivocan menos que aquellos en que
no, $P(\text{error} \mid a = 1) < P(\text{error} \mid a = 0)$, condición que se verifica con datos. Como usa
las decisiones y no las probabilidades, no depende de que los modelos estén calibrados, pero no informa
cuando los dos comparten un error.

El \textbf{consenso entre los vecinos} recuperados mide qué tan unánime es el voto de la
Ecuación~\ref{eq:marco_knn}, $\kappa(x) = \max_c \hat{p}_{\text{kNN}}(c \mid x)$. Vale 1 cuando los $k$
vecinos coinciden y depende de la representación y de la cobertura de la base: si la base no contiene casos
parecidos, los $k$ vecinos existen igual y no son informativos.

La calibración no crea información que la señal no contiene: un modelo que asigna 0.7 a todas sus
predicciones y acierta el 70\,\% de ellas está perfectamente calibrado, y su confianza no separa los casos
fáciles de los difíciles. Por eso una señal se evalúa con la curva riesgo--cobertura y no solo con una métrica
de calibración. Con una señal disponible, la decisión además deja de ser binaria.

\subsection{Tres salidas: afirmar, ofrecer diferenciales o derivar}
\label{sec:marco_tres_salidas}

\begingroup En la práctica clínica derivar no es la única alternativa a afirmar: un clínico que duda entre dos
o tres diagnósticos los enuncia como diagnóstico diferencial. Un sistema puede entonces tener tres salidas
(Tabla~\ref{tab:marco_salidas}): afirmar un diagnóstico, ofrecer diferenciales, es decir un conjunto de
candidatos, o derivar el caso cuando no se alcanza la confianza necesaria. Con dos umbrales
$\tau_1 > \tau_2$ sobre la señal, la política es\displaywidowpenalty=10000
\begin{equation}
    \pi(x) =
    \begin{cases}
        \text{afirmar} & \text{si } \kappa(x) \ge \tau_1,\\
        \text{diferencial} & \text{si } \tau_2 \le \kappa(x) < \tau_1,\\
        \text{derivar} & \text{si } \kappa(x) < \tau_2.
    \end{cases}
    \label{eq:marco_politica}
\end{equation}\endgroup

\begin{table}[H]
    \centering
    \caption{Salidas de un sistema con abstención: qué entrega cada una, cuándo corresponde y qué arriesga.}
    \label{tab:marco_salidas}
    \small
    \begin{tabularx}{\textwidth}{@{}>{\raggedright\arraybackslash}p{2.9cm}
        >{\raggedright\arraybackslash\hspace{0pt}}X >{\raggedright\arraybackslash\hspace{0pt}}X
        >{\raggedright\arraybackslash\hspace{0pt}}X@{}}
        \toprule
        \textbf{Salida} & \textbf{Qué entrega} & \textbf{Cuándo corresponde} & \textbf{Qué arriesga} \\
        \midrule
        Afirmar & Una categoría diagnóstica & Confianza alta
            & Un error que el lector no detecta por la forma del texto \\
        \filasep
        Ofrecer diferenciales & Un conjunto de categorías candidatas & Confianza intermedia
            & Informa menos, porque deja la elección entre candidatos al lector \\
        \filasep
        Derivar & Ninguna categoría; el caso pasa a un especialista & Confianza baja
            & No produce diagnóstico y traslada el caso a una consulta \\
        \bottomrule
    \end{tabularx}
\end{table}

La salida diferencial es una predicción por conjuntos, $S(x) \subseteq \{1, \dots, C\}$. El conjunto puede
formarse con las clases más probables de un modelo, hasta que su probabilidad acumulada alcance un umbral, o con
las predicciones de los distintos predictores cuando no coinciden. Cuando $|S(x)| = 1$, la salida coincide con
afirmar. Cuanto más caro es un error frente a derivar, más alto conviene fijar el umbral de afirmación, y esa
elección se hace en validación, nunca sobre el test.

Esta política decide cuánto afirma el sistema, pero no redacta nada. Cuando la salida tiene que ser un reporte
clínico, la decisión pasa a condicionar un texto, y un texto tiene problemas propios.
